\documentclass[11pt]{article}
\usepackage[utf8]{inputenc}
\usepackage[spanish]{babel}
\usepackage{amsmath,amssymb,amsthm}
\usepackage{graphicx}
\usepackage[margin=2.5cm]{geometry}
\usepackage{enumitem}
\title{Semana 5 -- Ejercicios}
\author{Tutorías -- Introducción al Cálculo}
\date{31/08--04/09/2026}
\begin{document}
\maketitle

\section{Ejercicios de la semana}

Resuelva las siguientes inecuaciones, indicando el conjunto solución como unión de intervalos.

\begin{enumerate}
    \item $|x-3| \le 5$. \hfill (Resp.\ $[-2,8]$)
    \item $|2x+1| > 3$. \hfill (Resp.\ $]-\infty,-2[\,\cup\,]1,\infty[$)
    \item $\dfrac{|x+1|}{|x-2|} \le 1$. \hfill (Resp.\ $\left]-\infty,\dfrac{1}{2}\right]$)
    \item $|3-x| \ge |x+1|$.
    \item $\sqrt{x+5} < x+1$.
    \item $\sqrt{2x-1} \le 3$. \hfill (Resp.\ $\left[\dfrac{1}{2},5\right]$)
    \item $\dfrac{(x^2-4)(x^2-9)}{(x+1)(x-5)} \ge 0$.
    \item $|x^2-4| < 5$.

    \item \textbf{(Nivel interrogación)} Resuelva $\sqrt{2x+3} < x$, determinando primero el
    dominio y justificando en qué paso es válido elevar al cuadrado.

    \item \textbf{(Nivel interrogación)} Resuelva $|x-1| \ge |2x+3|$ elevando al cuadrado ambos
    lados (justifique por qué esto es válido) y factorizando la diferencia de cuadrados
    resultante.

    \item \textbf{(Nivel interrogación)} Demuestre, a partir de la Definición 2.7 de valor
    absoluto (sin asumir propiedades ya demostradas en clase), que $\sqrt{x^2} = |x|$ para todo
    $x\in\mathbb{R}$.
\end{enumerate}

\end{document}
